\newglossaryentry{41viewmodel}{name={4+1 view model},description={Kruchten's organization of architecture into logical, process, development, and physical views, tied together by a set of scenarios (use cases).}}
\newglossaryentry{acceptancecriteria}{name={Acceptance criteria},description={Concrete conditions that must hold for a story to count as done, forming the seed of the tests that verify it.}}
\newglossaryentry{accidentalcomplexity}{name={Accidental complexity},description={Complexity that arises from our tools, choices, and mistakes, such as tangled dependencies, unclear names, and leaky abstractions, which good engineering can attack and reduce.}}
\newglossaryentry{actor}{name={Actor},description={Any entity outside the system that interacts with it to accomplish something, named by the role it plays rather than the person who plays it.}}
\newglossaryentry{actorgoallist}{name={Actor–goal list},description={A lightweight table with one row per goal tagged by its primary actor, serving as a one-page outline of the whole system.}}
\newglossaryentry{alternativeflow}{name={Alternative flow},description={A flow that attaches to a specific point in another flow, describing what happens when a stated condition interrupts the happy path.}}
\newglossaryentry{anchoring}{name={Anchoring},description={A cognitive bias in which the first number mentioned drags every subsequent estimate toward it, even when that number is arbitrary.}}
\newglossaryentry{appetite}{name={Appetite},description={A fixed time budget that constrains a feature's design, flipping the iron triangle so that time is fixed and scope flexes, with the deadline acting as a circuit breaker.}}
\newglossaryentry{architecturalpattern}{name={Architectural pattern},description={A reusable, named arrangement of modules and connectors that has repeatedly solved a recurring design problem.}}
\newglossaryentry{architecture}{name={Architecture},description={The set of significant, expensive-to-change design decisions about how a system is structured — its components, their responsibilities and interfaces, and the constraints among them — chosen to satisfy its most important quality requirements.}}
\newglossaryentry{architecturedecisionrecordsadrs}{name={Architecture Decision Records (ADRs)},description={Written records of significant, expensive-to-change design decisions and their rationale, which AI may draft but whose decision itself stays human-owned.}}
\newglossaryentry{architecturereview}{name={Architecture review},description={A structured examination of a system's design and component connections, conducted by people other than the designers to stress-test its assumptions and risks.}}
\newglossaryentry{attacktree}{name={Attack tree},description={A goal hierarchy for the adversary whose root is the attacker's goal, refined downward into concrete attacks whose leaves become defensive security requirements.}}
\newglossaryentry{attribute}{name={Attribute},description={A specific property of an entity that you care about, such as a module's size, a build's duration, or a defect's severity.}}
\newglossaryentry{automatedstaticanalysis}{name={Automated static analysis},description={The use of tools that examine source code or compiled artifacts without running them to report likely defects, style violations, or dangerous patterns.}}
\newglossaryentry{backlog}{name={Backlog},description={A prioritized list of the most important user stories or features in the users' language, begun in the proposal and later refined and estimated.}}
\newglossaryentry{basicflow}{name={Basic flow},description={The main success scenario or happy path: the shortest, cleanest numbered sequence of steps in which everything goes right and the goal is achieved.}}
\newglossaryentry{behaviordrivendevelopmentbdd}{name={Behavior-Driven Development (BDD)},description={An approach that writes requirements as concrete examples of behavior in customer-readable language, so the same scenarios can run as automated acceptance tests.}}
\newglossaryentry{blackboxtesting}{name={Black-box testing},description={Specification-based testing that chooses inputs from the specification while ignoring the code, so it finds missing logic and stays valid when the code is rewritten.}}
\newglossaryentry{boundaryvalueanalysis}{name={Boundary-value analysis},description={A black-box technique that targets the edges of equivalence classes by testing the value just below, at, and just above each boundary, where off-by-one defects cluster.}}
\newglossaryentry{boundedalternativeflow}{name={Bounded alternative flow},description={An alternative that applies across a range of steps rather than a single point, used for cross-cutting failures like timeouts or loss of connection.}}
\newglossaryentry{boxplot}{name={Boxplot},description={A box-and-whisker display of the five-number summary — minimum, Q1, median, Q3, and maximum — that marks outliers beyond the 1.5-IQR fences.}}
\newglossaryentry{brokenaccesscontrol}{name={Broken Access Control},description={The failure of the policy that users can only act within their intended permissions, letting users read, modify, or destroy data beyond their reach; defended by denying by default and deciding authorization on the server for every request.}}
\newglossaryentry{brokerpattern}{name={Broker pattern},description={A pattern that inserts an intermediary between clients and servers to locate, route, and return results, so the two sides need not know each other's location or identity.}}
\newglossaryentry{canarydeployment}{name={Canary deployment},description={Releasing a new version to a small slice of traffic first and watching its health metrics before progressively widening exposure to everyone.}}
\newglossaryentry{captheorem}{name={CAP theorem},description={The result that when a network partition splits a replicated data system, it can preserve at most two of consistency, availability, and partition tolerance.}}
\newglossaryentry{characterizationtest}{name={Characterization test},description={A test that documents what code actually does now by promoting observed outputs into assertions, using the running system itself as the oracle to pin down current behavior.}}
\newglossaryentry{ciatriad}{name={CIA triad},description={The three core properties security protects—confidentiality (only authorized parties can read data), integrity (data and code are not altered undetectably by unauthorized parties), and availability (the system is up when legitimate users need it).}}
\newglossaryentry{classdiagram}{name={Class diagram},description={The UML notation for the static structure of an object-oriented system: its classes, what each holds and does, and how they relate.}}
\newglossaryentry{clientserverpattern}{name={Client-server pattern},description={A pattern that designates a server as the keeper of a shared resource and clients as the many requesters that initiate contact and consume responses via a published protocol.}}
\newglossaryentry{cocomo}{name={COCOMO},description={Boehm's Constructive Cost Model, a family of power-law equations that estimate software effort in person-months from predicted size (KLOC or function points).}}
\newglossaryentry{codereview}{name={Code review},description={A lightweight practice in which at least one engineer other than the author reads and approves a change before it is merged, focused on intent and shared ownership.}}
\newglossaryentry{codingagents}{name={Coding agents},description={Generative-AI systems that can autonomously draft code, tests, requirements, and designs, and at higher rungs take an issue and open a pull request with little human typing.}}
\newglossaryentry{coefficientofdeterminationrsquared}{name={Coefficient of determination (R-squared)},description={The fraction of total variation in the response explained by the regression, computed as one minus the sum of squared errors over the total sum of squares.}}
\newglossaryentry{cohesion}{name={Cohesion},description={The degree to which a single module does one thing well; high cohesion is a property of a good decomposition.}}
\newglossaryentry{confidenceinterval}{name={Confidence interval},description={A range computed from a sample that, over many repeated samples, brackets the true population mean at a stated success rate such as 95\%.}}
\newglossaryentry{container}{name={Container},description={A package holding an application plus everything it needs to run, isolated from other software while sharing the host operating system's kernel.}}
\newglossaryentry{continuousdeployment}{name={Continuous deployment},description={A practice in which every change that passes the pipeline is deployed to production automatically, with no human in the loop, distinguished from delivery which only makes release possible on demand.}}
\newglossaryentry{continuousintegration}{name={continuous integration},description={A practice in which everyone merges into a shared mainline many times a day and an automated build-and-test pipeline verifies each merge.}}
\newglossaryentry{controlflowgraph}{name={Control-flow graph},description={A directed graph modeling a function, where nodes are basic blocks of straight-line statements and edges are the possible transfers of control between them.}}
\newglossaryentry{coupling}{name={Coupling},description={The degree to which modules depend on one another; low coupling, where modules interact only through their interfaces, is a property of a good decomposition.}}
\newglossaryentry{customerfounddefectcfd}{name={Customer-found defect (CFD)},description={A defect first reported by a customer in production, representing real pain that escaped every internal quality net.}}
\newglossaryentry{cyclomaticcomplexity}{name={Cyclomatic complexity},description={A metric equal to a function's decision points plus one (or E − N + 2 for a graph), measuring the number of independent paths and predicting testing effort.}}
\newglossaryentry{dataflowanalyzer}{name={Data-flow analyzer},description={A static tool that traces how values move through a program to find multi-line defects like use-before-assign, resource leaks, null dereferences, or tainted input.}}
\newglossaryentry{dataflownetwork}{name={Dataflow network},description={A generalization of the pipeline into a directed graph, where a filter may have several inputs and outputs and data can branch to parallel paths and merge back.}}
\newglossaryentry{decomposition}{name={Decomposition},description={The primary weapon against complexity: splitting a system into parts small enough to understand, with well-defined interfaces, so each can be reasoned about in isolation.}}
\newglossaryentry{deepdivereview}{name={Deep-dive review},description={A focused architecture review that examines one high-risk area exhaustively until the risk is either resolved or explicitly accepted.}}
\newglossaryentry{defectremovalefficiencydre}{name={Defect-removal efficiency (DRE)},description={The fraction of all defects ever present that were caught before release, computed as defects found before release divided by total defects before plus after release.}}
\newglossaryentry{defenseindepth}{name={Defense in depth},description={Placing several independent layers between an attacker and what matters—perimeter, hardened host, validating application, and encrypted data—chosen to fail independently so that defeating one does not defeat the next.}}
\newglossaryentry{definitionofdone}{name={Definition of Done},description={A shared, explicit checklist of what "finished" means for a backlog item, such as compiled, reviewed, tested, documented, and integrated.}}
\newglossaryentry{delphimethod}{name={Delphi method},description={A RAND-developed forecasting technique that combines expert anonymity, iteration, and statistical feedback to produce a group judgment reflecting knowledge rather than dominance.}}
\newglossaryentry{devops}{name={DevOps},description={The culture in which developing software and operating it are treated as one discipline, practiced by one team with shared tools and shared accountability.}}
\newglossaryentry{discoveryreview}{name={Discovery review},description={An early, broad architecture review that maps the problem space, names major components, and identifies big risks while the design is still forming.}}
\newglossaryentry{diseconomyofscale}{name={Diseconomy of scale},description={The property, captured by an exponent b greater than 1 in the power law Effort = a x Size\textsuperscript{b}, that bigger software costs disproportionately more per line due to communication and integration overhead.}}
\newglossaryentry{dorafourkeys}{name={DORA delivery metrics (the four keys)},description={The DORA outcome measures that jointly predict software-delivery performance—throughput (change lead time, deployment frequency, failed-deployment recovery time) and instability (change fail rate, deployment rework rate). Known as the ``four keys'' before deployment rework rate joined with the 2024 research.}}
\newglossaryentry{deploymentreworkrate}{name={Deployment rework rate},description={The share of deployments that are unplanned, shipped in response to a production incident—DORA's fifth delivery metric (2024), pairing with change fail rate on the instability side to make reactive work visible.}}
\newglossaryentry{dry}{name={DRY},description={The Don't Repeat Yourself rule that every piece of knowledge in a system should have a single, authoritative representation.}}
\newglossaryentry{elicitation}{name={Elicitation},description={The active work of drawing out needs that stakeholders cannot simply hand over, using questions, prompts, and situations that help people discover what they know but had not stated.}}
\newglossaryentry{empiricism}{name={Empiricism},description={The discipline's mature answer to changeable software: working like an experimental scientist by iterating in small steps, seeking fast feedback, and letting measured evidence rather than authority settle decisions.}}
\newglossaryentry{encapsulation}{name={Encapsulation},description={The language-enforced mechanism that keeps a module's private parts unreachable from outside, so clients cannot accidentally depend on its internals.}}
\newglossaryentry{essentialcomplexity}{name={Essential complexity},description={Complexity inherent in the problem itself, such as a tax program being complicated because tax law is complicated, which no cleverness can remove.}}
\newglossaryentry{exploit}{name={Exploit},description={The concrete technique or code that turns a specific vulnerability into an actual compromise.}}
\newglossaryentry{extend}{name={Extend},description={A relationship that adds optional behavior onto a base use case at an extension point, running only when its condition holds, with the base oblivious to its extender.}}
\newglossaryentry{extensionpoint}{name={Extension point},description={A named location in a flow where optional behavior may be attached, letting an extending use case plug in without cluttering the base flow.}}
\newglossaryentry{extremeprogrammingxp}{name={Extreme Programming (XP)},description={An agile approach that prescribes a team's engineering discipline, turning known-good practices like testing, review, and integration up to the extreme.}}
\newglossaryentry{faganinspection}{name={Fagan inspection},description={The archetypal formal inspection pioneered by Michael Fagan, defined by a small group, a fixed sequence of phases, assigned roles, and measurement of the process itself.}}
\newglossaryentry{failure}{name={Failure},description={The observable event when the running system deviates from correct behavior, such as a report printing the wrong total.}}
\newglossaryentry{falsenegative}{name={False negative},description={A real defect that the tool fails to report, giving false reassurance that broken code is clean.}}
\newglossaryentry{falsepositive}{name={False positive},description={A tool report of a defect where the code is actually fine, an alarm the analyzer raises because it cannot prove the code correct.}}
\newglossaryentry{fault}{name={Fault},description={The resulting flaw in the software (the defect in the code), such as an off-by-one in a loop, that stays latent until some execution triggers it.}}
\newglossaryentry{featureflag}{name={Feature flag},description={A runtime conditional that turns a code path on or off without redeploying, decoupling the act of deploying code from the act of releasing a feature.}}
\newglossaryentry{functionalrequirement}{name={Functional requirement},description={A requirement describing a behavior the system does in response to input, usually testable with a concrete given-input, expected-output example.}}
\newglossaryentry{gherkin}{name={Gherkin},description={A structured-English Given/When/Then format that expresses acceptance criteria as concrete scenarios a customer can read.}}
\newglossaryentry{goal}{name={Goal},description={What the primary actor wants the interaction to achieve, stated as a result in the world rather than as a UI action or software feature.}}
\newglossaryentry{goalquestionmetricgqm}{name={Goal-Question-Metric (GQM)},description={A top-down measurement method that derives metrics from concrete questions, which are in turn derived from an explicitly stated business or engineering goal.}}
\newglossaryentry{goodhartslaw}{name={Goodhart's Law},description={The principle that once a metric becomes a target tied to incentives, people optimize the metric itself rather than the goal it was meant to represent.}}
\newglossaryentry{growitculture}{name={grow-it culture},description={A development philosophy that starts with something small that works and evolves it through many small, reality-validated changes, discovering requirements through working software.}}
\newglossaryentry{hallucinatedrequirements}{name={Hallucinated requirements},description={Plausible-sounding requirements an LLM invents that are actually wrong, whose fluency can manufacture a false sense of consensus about real user needs.}}
\newglossaryentry{hillchart}{name={Hill chart},description={A progress display, popularized by Shape Up, that shows how much of the work is still uncertain (uphill) versus mere execution (downhill), rather than just quantity remaining.}}
\newglossaryentry{histogram}{name={Histogram},description={A display of the shape of a single continuous variable that divides its range into equal bins and draws a bar whose height is the count of values in each bin.}}
\newglossaryentry{horizontalscaling}{name={Horizontal scaling},description={Meeting rising load by running more interchangeable copies of a service behind a load balancer rather than by buying a single bigger machine.}}
\newglossaryentry{humbleview}{name={Humble view},description={The discipline of keeping a view as logic-free as possible so that business rules live in a model or view-model where automated tests can reach them.}}
\newglossaryentry{include}{name={Include},description={A relationship that factors shared, unconditional behavior into its own use case that other use cases invoke every time they reach that point.}}
\newglossaryentry{informationhiding}{name={Information hiding},description={Designing each module around a decision likely to change and hiding that decision behind an interface, so the rest of the system depends on the interface rather than the decision.}}
\newglossaryentry{injection}{name={Injection},description={What happens when untrusted input is handed to an interpreter and part of that input is executed as a command rather than treated as data, with SQL injection and cross-site scripting as the canonical cases.}}
\newglossaryentry{insecuredesign}{name={Insecure Design},description={Missing or ineffective control design—a security weakness that a perfect implementation would faithfully preserve, catchable only by threat modeling and design review because there is no defect on any single line to find.}}
\newglossaryentry{integrationtesting}{name={Integration testing},description={Testing that exercises the interfaces and interactions between units, surfacing mismatched assumptions, wrong data formats, and protocol errors.}}
\newglossaryentry{invest}{name={INVEST},description={A checklist for judging whether a story is ready to build: Independent, Negotiable, Valuable, Estimable, Small, and Testable.}}
\newglossaryentry{investedexpertreviewer}{name={Invested expert reviewer},description={A reviewer who knows a part of the system deeply and has a stake in its long-term health, bringing context that is not visible in the diff.}}
\newglossaryentry{irontriangle}{name={Iron triangle},description={The set of three coupled project constraints (scope, cost, and time) that cannot all be fixed independently, forcing deliberate trade-offs.}}
\newglossaryentry{kanomodel}{name={Kano model},description={A mental model that categorizes features as must-be, performance, or delighter by how user satisfaction responds to their presence or absence.}}
\newglossaryentry{ladderofassistance}{name={Ladder of assistance},description={The progression of AI coding capability from autocomplete through chat assistant, in-IDE agent, and autonomous agent up to a multi-agent swarm, along which the human's job shifts from typing to specifying, supervising, and verifying.}}
\newglossaryentry{latentneed}{name={Latent need},description={A need the user does not know they want until they see it, but which delights them when they do, and where competitive advantage lives.}}
\newglossaryentry{layeredpattern}{name={Layered pattern},description={A structure that stacks a system into horizontal tiers, where each layer offers services to the layer above and consumes services from the layer below, with dependencies pointing downward only.}}
\newglossaryentry{leastprivilege}{name={Least privilege},description={The habit of giving every user, service, and token the minimum access it needs and no more, so that a compromise of one part cannot reach the whole.}}
\newglossaryentry{legacycode}{name={Legacy code},description={Code without tests—code you are afraid to change—whose actual behavior is pinned down nowhere except in the code itself, regardless of its age.}}
\newglossaryentry{linter}{name={Linter},description={A static tool that flags stylistic issues, suspicious constructs, and small correctness hazards such as unused variables or always-true comparisons.}}
\newglossaryentry{mapreduce}{name={MapReduce},description={A big-data dataflow style that expresses a computation as a map step transforming each record independently, followed by a reduce step aggregating the mapped results by key.}}
\newglossaryentry{metric}{name={Metric},description={A rule that assigns a number or symbol to an attribute in a way that preserves something true about the world, so entities differing in the attribute differ in the metric.}}
\newglossaryentry{modelviewcontrollermvc}{name={Model-View-Controller (MVC)},description={A pattern that splits an interactive application into a model (data and rules), a view (rendering), and a controller (input handling) so each concern can change and be tested independently.}}
\newglossaryentry{modifiedconditiondecisioncoveragemcdc}{name={Modified Condition/Decision Coverage (MC/DC)},description={A strong criterion requiring each condition in a compound decision to be shown, via a pair of tests, to independently flip the decision's outcome.}}
\newglossaryentry{modularity}{name={Modularity},description={The practice of building a system out of parts that can each be understood, built, tested, and changed largely on their own.}}
\newglossaryentry{module}{name={Module},description={A unit of software with a boundary — a set of things it provides to the rest of the system and a set of things it keeps to itself — whose value comes almost entirely from what it hides.}}
\newglossaryentry{moscow}{name={MoSCoW},description={A DSDM prioritization scheme that sorts requirements into Must have, Should have, Could have, and Won't have (this time) buckets.}}
\newglossaryentry{multitenancy}{name={Multi-tenancy},description={A design in which one running instance of the system serves many customers (tenants) whose data is kept logically separate inside shared infrastructure.}}
\newglossaryentry{multiplicity}{name={Multiplicity},description={An annotation on the ends of a class-diagram line stating how many objects participate, such as 1, 0..1, 1..*, or 0..*.}}
\newglossaryentry{mutationtesting}{name={Mutation testing},description={Grading a test suite by planting small deliberate defects called mutants and measuring the fraction the suite kills (detects) versus lets survive.}}
\newglossaryentry{observerpattern}{name={Observer pattern},description={A pattern in which interested parties subscribe to a subject that notifies its subscribers of state changes without knowing their concrete types.}}
\newglossaryentry{oracleproblem}{name={Oracle problem},description={The difficulty that a generated test only encodes what the model thinks correct behavior is, so coverage can look great while both code and test share the same misunderstanding of the specification.}}
\newglossaryentry{ordinaryleastsquaresols}{name={Ordinary least squares (OLS)},description={A regression method that fits the line minimizing the sum of squared vertical residuals between observed and predicted values.}}
\newglossaryentry{outcomeengineeringo16gmanifesto}{name={Outcome Engineering (o16g manifesto)},description={Cory Ondrejka's 2026 position paper arguing that agentic development should be graded on verified outcomes rather than code, framed around four shifts—creation not code, cost not time, capacity not backlog, certainty not vibes—and sixteen principles.}}
\newglossaryentry{owasptop10}{name={OWASP Top 10},description={A widely used awareness document that ranks the ten most critical categories of web-application security risk, assembled from empirical vulnerability data plus a practitioner survey and revised roughly every four years.}}
\newglossaryentry{parameterizedquery}{name={Parameterized query},description={A query in which the SQL text is fixed and compiled first and untrusted input is bound in afterward as pure data via a placeholder, so it can never change the query's structure.}}
\newglossaryentry{pipesandfilters}{name={Pipes and filters},description={A dataflow pattern that arranges computation as a linear chain where each filter reads from an input pipe, transforms the data, and writes to an output pipe feeding the next filter.}}
\newglossaryentry{plandrivenculture}{name={plan-driven culture},description={A development philosophy that decides in full detail what to build before building it, freezing requirements and prizing predictability, documentation, and up-front analysis.}}
\newglossaryentry{planningpoker}{name={Planning Poker},description={A fast agile form of Wideband Delphi in which estimators privately pick point cards, reveal simultaneously, and let the high and low estimators explain before re-voting until they converge.}}
\newglossaryentry{precision}{name={Precision},description={The fraction of a tool's alarms that correspond to real defects; low precision means many false positives.}}
\newglossaryentry{primaryactor}{name={Primary actor},description={The actor who initiates a use case to reach a goal of their own, whose desire is the reason the use case exists.}}
\newglossaryentry{productivityparadox}{name={Productivity paradox},description={The gap between perceived and measured productivity, where developers feel faster with AI yet controlled studies can show them slower because reviewing and integrating generated output dominates.}}
\newglossaryentry{propertybasedtesting}{name={Property-based testing},description={Testing in which a tool generates many random inputs and checks that a stated property or invariant holds on each, rather than hard-coding one expected answer.}}
\newglossaryentry{proposal}{name={Proposal},description={The first deliverable, due around week 3, in which a team commits to a real problem, a specific user, a plausible scope, an intended architecture, a process, and its biggest risks.}}
\newglossaryentry{psychologicalsafety}{name={Psychological safety},description={The shared belief that a team member can raise a problem, admit a mistake, or ask a basic question without being punished or embarrassed, identified as the research-backed foundation of team performance.}}
\newglossaryentry{publishsubscribe}{name={Publish–subscribe},description={The network-scale form of observer, where publishers emit events to named topics and subscribers register interest in topics rather than in a specific subject.}}
\newglossaryentry{qualityrequirement}{name={Quality requirement},description={A non-functional requirement describing a property the system must have while behaving (speed, security, availability), which needs a measurable fit criterion to make "done" a fact.}}
\newglossaryentry{qualityattributescenario}{name={Quality-attribute scenario},description={A short, structured statement — naming stimulus, environment, response, and response measure — that turns a vague quality requirement into an observable event with a measurable outcome.}}
\newglossaryentry{recall}{name={Recall},description={The fraction of real defects a tool catches; low recall means many false negatives.}}
\newglossaryentry{refactoring}{name={refactoring},description={Continuously improving the design of existing code without changing its behavior, kept safe by an automated test suite.}}
\newglossaryentry{referencestory}{name={Reference story},description={A small, well-understood item the whole team agrees to assign a fixed point value, against which all other estimates are calibrated as bigger or smaller.}}
\newglossaryentry{requirement}{name={Requirement},description={A statement of something the software must do or a quality it must have, elicited from diverse stakeholders who often describe solutions rather than their underlying problem.}}
\newglossaryentry{requirementsanalysis}{name={Requirements analysis},description={The discipline of sizing work, estimating effort, and ranking requirements by value, cost, and risk to turn a backlog of wishes into a defensible plan.}}
\newglossaryentry{requirementstraceability}{name={Requirements traceability},description={The ability to follow a requirement forward to the backlog items, code, and tests that realize it, and backward from any of those to the requirement that justifies it.}}
\newglossaryentry{restrepresentationalstatetransfer}{name={REST (Representational State Transfer)},description={A set of conventions—resources named by URIs, a fixed set of HTTP verbs, exchanged representations, and statelessness—that make the client-server pattern scale across the web.}}
\newglossaryentry{retrospectivereview}{name={Retrospective review},description={An architecture review conducted after the system runs to learn which design assumptions held and where reality diverged from the plan.}}
\newglossaryentry{scalesofmeasurement}{name={Scales of measurement},description={The four classic levels of data — nominal, ordinal, interval, and ratio — that determine which statistics and charts are legal for a given metric.}}
\newglossaryentry{scenario}{name={Scenario},description={A concrete story of a specific person using the system to accomplish a goal, step by step in a specific situation, exposing gaps and branches that terse formats hide.}}
\newglossaryentry{scopehammering}{name={Scope hammering},description={The Shape Up discipline of continually attacking scope under a fixed appetite, actively looking for use cases and implementation pieces that can be cut.}}
\newglossaryentry{scrum}{name={Scrum},description={A lightweight framework that organizes a team's time and decisions into short cycles delivering working software, without prescribing how to write, design, or test code.}}
\newglossaryentry{severity}{name={Severity},description={An ordinal classification of a defect by the impact of the failure it causes, ranging from critical (S1) down to low (S4).}}
\newglossaryentry{shapeup}{name={Shape Up},description={A method that fixes time and flexes scope, keeps no backlog, and hands a team a whole shaped problem rather than a list of tasks.}}
\newglossaryentry{shareddatarepositorypattern}{name={Shared-data (repository) pattern},description={A design in which independent accessor components read and write a single central data store that serves as the authoritative source of truth.}}
\newglossaryentry{shiftingleft}{name={Shifting left},description={Moving security activity toward the start of the lifecycle—threat modeling during requirements, secure design review before code, and scanning on every commit—because the earlier a vulnerability is caught, the cheaper it is to fix.}}
\newglossaryentry{skeletalsystem}{name={Skeletal system},description={The first status-report milestone: a thin but fully wired system, built as a walking skeleton, that forces integration risk to surface early rather than at the end of the term.}}
\newglossaryentry{slsa}{name={SLSA},description={Supply-chain Levels for Software Artifacts, a framework of build-integrity levels from L0 (no guarantees) through L3 (a hardened build platform with strong tamper resistance) that grades how trustworthy an artifact's origin story is.}}
\newglossaryentry{smoketest}{name={Smoke test},description={A fast, shallow pass that answers whether a build is fundamentally alive, run as a gate immediately after a build or deployment.}}
\newglossaryentry{softwareasaservicesaas}{name={Software as a service (SaaS)},description={A delivery model in which the software runs on machines the vendor controls and users reach it over the network, so nothing ships to the customer.}}
\newglossaryentry{softwarebillofmaterialssbom}{name={Software bill of materials (SBOM)},description={A complete inventory of every component in a release, which turns "patch the dependency you did not know you had" from impossible into a database query; dominant standards are CycloneDX and SPDX.}}
\newglossaryentry{softwarecrisis}{name={Software crisis},description={The name given at the 1968 NATO conference to the pattern of large software projects running chronically late, over budget, and failing outright, which motivated developing software as an engineering discipline.}}
\newglossaryentry{softwaredevelopmentprocess}{name={software development process},description={A repeatable way of organizing the activities that turn an idea into deployed, maintained software, serving as a team's standing answer to complexity, change, defects, and coordination.}}
\newglossaryentry{softwareengineering}{name={Software engineering},description={The disciplined application of principles, methods, and tools to build and evolve software that is useful, correct enough, and economical to change, usually by a team under changing requirements and at a scale that defeats ad-hoc effort.}}
\newglossaryentry{softwareinspection}{name={Software inspection},description={The most formal static-checking technique: a disciplined, role-based, line-by-line examination of a work product for the explicit purpose of finding defects.}}
\newglossaryentry{softwareproductlineengineering}{name={Software product-line engineering},description={Building a family of related products from a shared set of assets in a planned way, separating commonalities from designed-in variabilities.}}
\newglossaryentry{solid}{name={SOLID},description={A bundle of five object-oriented design principles — single responsibility, open-closed, Liskov substitution, interface segregation, and dependency inversion.}}
\newglossaryentry{specificalternativeflow}{name={Specific alternative flow},description={An alternative that branches at one identified step, handles a particular condition, and then either rejoins the basic flow or ends the use case.}}
\newglossaryentry{spiralmodel}{name={spiral model},description={An explicitly risk-driven process framework that proceeds in outward loops, each attacking the biggest remaining risk before developing and verifying a piece.}}
\newglossaryentry{sprint}{name={sprint},description={A fixed-length Scrum iteration, usually one to four weeks, that produces a potentially shippable increment of the product.}}
\newglossaryentry{standarddeviation}{name={Standard deviation},description={The square root of the variance, expressing in original units the typical distance of a data value from the sample mean.}}
\newglossaryentry{staticchecking}{name={Static checking},description={The discipline of finding faults by examining a program's design documents, source code, and types rather than by executing it.}}
\newglossaryentry{statusreport}{name={Status report},description={A milestone deliverable that demonstrates a running system and gives an honest account of what is done, what changed, and what is blocked, filed once for the skeletal system and once for the viable system.}}
\newglossaryentry{storypoint}{name={Story point},description={A unit of relative size (not time) expressing how much effort, complexity, and uncertainty a story carries compared to a reference story.}}
\newglossaryentry{storyboard}{name={Storyboard},description={A scenario told in pictures as a short sequence of rough, low-fidelity sketches, one frame per step, showing what the user sees and does at each moment.}}
\newglossaryentry{stranglerfig}{name={Strangler fig},description={A pattern for incrementally replacing a legacy system by routing one capability at a time to a new implementation behind an interception facade until the old code is retired.}}
\newglossaryentry{subflow}{name={Subflow},description={A named, reusable chunk of steps within a single use case that the main flow can invoke by name, like a local helper function.}}
\newglossaryentry{supportingactor}{name={Supporting actor},description={An external system or person the system calls on to get the job done, but whose own goals are not the point of the use case.}}
\newglossaryentry{swebench}{name={SWE-bench},description={A benchmark that tasks agents with fixing a real GitHub issue end to end, used to measure automated program-repair progress but whose passing tests are evidence, not proof, of a correct fix.}}
\newglossaryentry{swtbench}{name={SWT-bench},description={A benchmark that tasks agents with generating a bug-reproducing test, used to track progress in automated test generation.}}
\newglossaryentry{technicaldebt}{name={Technical debt},description={The metaphor for the future cost of a shortcut taken today, whose interest is that every later change to that code costs more than it otherwise would.}}
\newglossaryentry{testdouble}{name={Test double},description={A stand-in for a real dependency: a stub returns canned answers, a mock also records how it was called, and a fake is a lightweight working implementation.}}
\newglossaryentry{testoracle}{name={Test oracle},description={Whatever decides, for a given input, whether the actual output is correct, ranging from an exact expected value down to human judgment.}}
\newglossaryentry{testadequacycriterion}{name={Test-adequacy criterion},description={A checkable coverage condition that tells you whether a suite is done with respect to that criterion, turning "have we tested enough?" into a countable question.}}
\newglossaryentry{testdrivendevelopmenttdd}{name={test-driven development (TDD)},description={A tight red-green-refactor loop in which you write a failing test, add the least code to make it pass, then improve the design while the tests stay green.}}
\newglossaryentry{testing}{name={Testing},description={Running the software on chosen inputs and checking the outputs against expectations; it can show the presence of defects but never prove their absence.}}
\newglossaryentry{testingpyramid}{name={Testing pyramid},description={A strategy of writing many fast unit tests, fewer integration tests, and only a handful of slow end-to-end tests, driven by the cost gradient between levels.}}
\newglossaryentry{threat}{name={Threat},description={A potential source of harm—the actor or event that might act against you, such as a criminal crew, a disgruntled insider, or an automated bot sweeping the internet.}}
\newglossaryentry{trunkbaseddevelopment}{name={Trunk-based development},description={A branching policy in which the whole team commits to a single shared mainline, directly or through branches that live only hours to a couple of days.}}
\newglossaryentry{typechecker}{name={Type checker},description={A static tool that verifies operations are applied to compatible kinds of data, rejecting whole classes of defect before the program runs.}}
\newglossaryentry{unboundedstream}{name={Unbounded stream},description={Data that never ends, forcing computations to produce results incrementally using windows, out-of-order and late-data handling, and backpressure.}}
\newglossaryentry{unittest}{name={Unit test},description={A test that exercises the smallest independently testable piece, typically one function or class, in isolation from the rest of the system.}}
\newglossaryentry{usecase}{name={Use case},description={A complete narrative of one goal's interaction between a user and the system, from opening move to success or controlled failure, including everything that can go wrong.}}
\newglossaryentry{usecasediagram}{name={Use-case diagram},description={A one-glance map of who the actors are, what goals the system offers, and how those goals relate, showing scope and relationships rather than flows.}}
\newglossaryentry{userstory}{name={user story},description={A short description of a feature told from the user's point of view, treated as a promise to have a conversation rather than a full specification.}}
\newglossaryentry{vmodel}{name={V-model},description={A refinement of waterfall that pairs each level of design with a matching level of testing, from unit up through integration, system, and acceptance.}}
\newglossaryentry{velocity}{name={Velocity},description={The number of story points a team actually completes in one sprint, measured (not decreed) and averaged to forecast remaining work.}}
\newglossaryentry{viablesystem}{name={Viable system},description={A minimum viable product that a real user could actually use to accomplish the one or two core tasks correctly, even if it lacks polish, secondary features, and edge-case handling.}}
\newglossaryentry{view}{name={View},description={A representation of the system drawn from one stakeholder perspective, since no single diagram can capture the whole.}}
\newglossaryentry{vulnerability}{name={Vulnerability},description={A weakness in the system, such as a missing bounds check, an unescaped query, an over-privileged account, or a dependency with a known flaw, that you own and can retire before any exploit exists.}}
\newglossaryentry{vulnerabilitybasedtrust}{name={Vulnerability-based trust},description={The trust a team builds by openly admitting what its members do not know, produced through habits like leaders confessing mistakes first and thanking people for surfacing bad news early.}}
\newglossaryentry{walkingskeleton}{name={Walking skeleton},description={A tiny end-to-end implementation that sends a real request through every architectural layer to the database and back, proving that every connection works even though it does almost nothing useful yet.}}
\newglossaryentry{waterfallmodel}{name={waterfall model},description={A process that runs a project as a single pass through a strict sequence of phases, each completed and signed off before the next begins.}}
\newglossaryentry{weaklinks}{name={Weak links},description={Structural weaknesses in how open-source packages are maintained—such as under-supported solo maintainers or account-recovery email running through expired domains—so that compromising a few high-profile developer accounts could poison much of an ecosystem.}}
\newglossaryentry{whiteboxtesting}{name={White-box testing},description={Structure-based testing that chooses inputs by looking at the code to exercise its statements, branches, and paths that specification-based tests never reach.}}
\newglossaryentry{widebanddelphi}{name={Wideband Delphi},description={Boehm's structured estimation cycle in which experts estimate privately, reveal together, have outliers explain their reasoning, and then re-estimate.}}
\newglossaryentry{wsjfweightedshortestjobfirst}{name={WSJF (Weighted Shortest Job First)},description={A prioritization method that ranks work by cost of delay divided by job size, a member of the same value-over-cost family used throughout the chapter.}}
